\DocumentMetadata{
  pdfversion=2.0,
  pdfstandard=ua-2,
  lang=en-US,
  tagging=on,
  tagging-setup={math/setup=mathml-SE,tabsorder=structure}
}
\documentclass[11pt,letterpaper]{article}
\usepackage[margin=0.68in,includefoot]{geometry}
\usepackage{newcomputermodern}
\usepackage{amsmath,amscd,mathtools}
\usepackage{tikz,adjustbox}
\providecommand{\square}{\mdlgwhtsquare}
\usepackage[hidelinks]{hyperref}
\hypersetup{
  pdftitle={Analysis PhD Exam, May 2024},
  pdfauthor={University of Florida Department of Mathematics},
  pdfsubject={Graduate mathematics examination},
  pdfdisplaydoctitle=true,
  bookmarksopen=true
}
\setcounter{secnumdepth}{0}
\setlength{\parindent}{0pt}
\setlength{\parskip}{0.28em}
\setlength{\emergencystretch}{3em}
\setlength{\leftmargini}{2.15em}
\setlength{\leftmarginii}{2.65em}
\setlength{\leftmarginiii}{3em}
\pagestyle{empty}
\raggedbottom
\allowdisplaybreaks
\NewTaggingSocketPlug{block/list/label}{examproblem}{%
  \tagstructbegin{tag=Lbl}\tagstructend%
  \tagstructbegin{tag=\LBody}%
  \tagstructbegin{tag=\ProblemHeadingTag,title-o={\ProblemHeadingText}}%
  \tagmcbegin{tag=\ProblemHeadingTag,actualtext={\ProblemHeadingText}}%
  #2%
  \tagmcend%
  \tagstructend%
}
\newenvironment{examproblems}[2]{%
  \def\ProblemHeadingTag{#2}%
  \AssignTaggingSocketPlug{block/list/label}{examproblem}%
  \begin{enumerate}%
  \setcounter{enumi}{#1}%
  \renewcommand{\labelenumi}{\textbf{\theenumi.}}%
  \setlength{\topsep}{0.35em}%
  \setlength{\partopsep}{0pt}%
  \setlength{\itemsep}{0.48em}%
  \setlength{\parsep}{0.18em}%
}{\end{enumerate}}
\newenvironment{examsubparts}{%
  \AssignTaggingSocketPlug{block/list/label}{default}%
  \begin{enumerate}%
  \setlength{\topsep}{0.18em}%
  \setlength{\partopsep}{0pt}%
  \setlength{\itemsep}{0.16em}%
  \setlength{\parsep}{0.1em}%
}{\end{enumerate}}
\newenvironment{examinstructions}{%
  \par\begingroup\itshape
}{%
  \par\endgroup\vspace{0.18em}
}
\makeatletter
\renewcommand\subsection{\@startsection{subsection}{2}{0pt}%
  {-1.25ex plus -.25ex minus -.1ex}%
  {0.55ex plus .1ex}%
  {\normalfont\large\bfseries}}
\renewcommand{\@maketitle}{%
  \newpage\null\vskip -1em%
  {\footnotesize\bfseries
    \href{https://www.ufl.edu/}{University of Florida}, \href{https://math.ufl.edu/}{Department of Mathematics}\par}%
  \vskip -0.5em%
  \tagstructbegin{tag=H1,title={Analysis PhD Exam, May 2024}}%
  \tagpdfparaOff%
  \tagmcbegin{tag=H1}%
    {\LARGE\bfseries\href{https://gma.math.ufl.edu/exams/analysis-phd/}{Analysis PhD Exam}, May 2024\par}%
  \tagmcend%
  \tagpdfparaOn%
  \tagstructend%
  \vskip 0.25em%
  \tagmcbegin{artifact}%
    \hrule height 1pt%
  \tagmcend%
  \vskip 1em%
}
\makeatother
\title{Analysis PhD Exam, May 2024}
\author{}
\date{}
\begin{document}
\maketitle
\thispagestyle{empty}
\begin{examinstructions}
Write solutions in a neat and logical fashion, giving complete reasons for all steps.
\end{examinstructions}
\begin{examinstructions}
Attempt SIX problems.
\end{examinstructions}
\begin{examproblems}{0}{H2}
\def\ProblemHeadingText{Problem 1.}
\item
Consider Lebesgue measure on the real line~\(\mathbb{R}\). Give an example for each of the following, if possible. If not possible, give a brief explanation. A sequence \(f_n\) in \(L^1(\mathbb{R})\) converging to an~\(f\) in \(L^1(\mathbb{R})\)...
\begin{examsubparts}
\item[\textnormal{a)}]
...in the \(L^1\) norm but not in measure,
\item[\textnormal{b)}]
...in measure but not in the \(L^1\) norm,
\item[\textnormal{c)}]
...in the \(L^1\) norm but not a.e.,
\item[\textnormal{d)}]
...a.e. but not in measure.
\end{examsubparts}
\def\ProblemHeadingText{Problem 2.}
\item
State Tonelli’s theorem, and give an example to show that the~\(\sigma\)-finiteness hypothesis is necessary.
\def\ProblemHeadingText{Problem 3.}
\item
For~\(\sigma\)-finite measures~\(\mu\) and~\(\nu\), define what it means for~\(\mu\) to be absolutely continuous with respect to~\(\nu\), and singular with respect to~\(\nu\). State the Lebesgue–Radon–Nikodym theorem.
\def\ProblemHeadingText{Problem 4.}
\item
Let \((X,\mathcal{M},\mu)\) be a~\(\sigma\)-finite measure space. Suppose that~\(g\) is a measurable function with the property that \(gf\in L^2\) for all \(f\in L^2\). Prove that \(g\in L^\infty\).
\def\ProblemHeadingText{Problem 5.}
\item
Fix \(1<p<\infty\), let~\(\mu\) be a~\(\sigma\)-finite measure, and let \((f_n)\) be a sequence in \(L^p(\mu)\). Suppose that there exists a function \(f\in L^p(\mu)\) such that
\[
\int f_n g\,d\mu \to \int fg\,d\mu
\]
 for every \(g\in L^q(\mu)\), where \(\frac{1}{p}+\frac{1}{q}=1\).
\begin{examsubparts}
\item[\textnormal{a)}]
Prove that \(\sup_n\{\lVert f_n\rVert_p\}<+\infty\).
\item[\textnormal{b)}]
Prove that \(\lVert f\rVert_p\leq\limsup\lVert f_n\rVert_p\).
\item[\textnormal{c)}]
Give an example to show that strict inequality can hold in part (b).
\end{examsubparts}
\def\ProblemHeadingText{Problem 6.}
\item
Let~\(\mathcal{X}\) be a normed vector space. Say \(x_n\to x\) weakly if \(f(x_n)\to f(x)\) for all \(f\in\mathcal{X}^*\). Prove that if~\(\mathcal{M}\) is a norm closed subspace of~\(\mathcal{X}\), and \((x_n)\) is a sequence in~\(\mathcal{M}\) converging weakly to~\(x\), then \(x\in\mathcal{M}\).
\def\ProblemHeadingText{Problem 7.}
\item
Let~\(\mu\) be a finite signed Borel measure on \([0,1]\). Suppose that
\[
\int_0^1 \sin(2\pi nx)\,d\mu(x)=0 \quad\text{and}\quad \int_0^1 \cos(2\pi nx)\,d\mu(x)=0
\]
 for all \(n=0,1,2,\ldots\). Prove that \(\mu=0\). Does it suffice to assume only the vanishing of the sine integrals? Prove, or give a counterexample.
\def\ProblemHeadingText{Problem 8.}
\item
Short answer.
\begin{examsubparts}
\item[\textnormal{a)}]
Give a brief explanation of how the Fourier transform is defined on \(L^2(\mathbb{R})\).
\item[\textnormal{b)}]
Sketch a proof that the Banach space \(L^1(\mathbb{R})\) is not reflexive.
\item[\textnormal{c)}]
Explain why the norm in \(\ell^1(\mathbb{N})\) cannot be given by an inner product.
\end{examsubparts}
\end{examproblems}
\end{document}
